% !PRINTED 545
\weberpage{545}
Es ist aber
\[
\begin{aligned}
K(-3) &= \tfrac{1}{3}, & K(-84) &= 4, & K(-20) &= 2, & K(-35) &= 2, \\
K(+140) &= 4, & K(5) &= 1, & K(21) &= 2, & K(12) &= 2^{1)},
\end{aligned}
\]
\begingroup\renewcommand{\thefootnote}{1)}\footnotetext[1]{Nach \gesperrt{Gauss}scher Bezeichnung sind die Klassenzahlen zweiter Art zu nehmen.}\endgroup
wie man aus den \gesperrt{Gauss}schen Tafeln oder hier auch leicht durch direkte Abzählung findet.

Ferner ist nach den \gesperrt{Legendre}schen oder \gesperrt{Degen}schen Tafeln:
\[
\begin{aligned}
\log\left(\frac{T+\sqrt{140}\,U}{2}\right) &= \log(6+\sqrt{35}) = \log\frac{(\sqrt{5}+\sqrt{7})^2}{2}, \\
\log\left(\frac{T+\sqrt{5}\,U}{2}\right) &= \log\left(\frac{1+\sqrt{5}}{2}\right)^2 = \tfrac{2}{3}\log(2+\sqrt{5}), \\
\log\frac{T+\sqrt{21}\,U}{2} &= \log\frac{5+\sqrt{21}}{2} = \tfrac{1}{3}\log(55+12\sqrt{21}) \\
&= \log\left(\frac{\sqrt{3}+\sqrt{7}}{2}\right)^2, \\
\log\left(\frac{T+\sqrt{12}\,U}{2}\right) &= \log(2+\sqrt{3}),
\end{aligned}
\]
und daraus erhält man:
\[
\left(\frac{f(\sqrt{-105})}{\sqrt[4]{2}}\right)^{12} = (2+\sqrt{5})^2(55+12\sqrt{21})(6+\sqrt{35})(2+\sqrt{3})^3,
\]
oder was damit gleichbedeutend ist:
\[
\sqrt{2}^{13}f(\sqrt{-105})^6 = (1+\sqrt{5})^3(\sqrt{3}+\sqrt{7})^3(\sqrt{5}+\sqrt{7})(1+\sqrt{3})^3.
\]

\begin{center}
\textbf{\S\,144. Berechnung einiger weiterer Klasseninvarianten.}
\end{center}

Nächst den Diskriminanten, bei denen in jedem Geschlecht nur \gesperrt{eine} Formenklasse vorkommt, die wir im vorhergehenden Paragraphen betrachtet haben, geben die einfachsten Resultate die, welche in jedem Geschlecht zwei Formenklassen enthalten, und unter diesen wieder die, bei denen nur zwei Geschlechter vorkommen. Die Klasseninvarianten für diese Diskriminanten $-4m$ sind die Wurzeln quadratischer Gleichungen, deren Koeffizienten nur eine Quadratwurzel enthalten. Wir setzen wieder Stammdiskriminanten voraus, und erhalten nach der \gesperrt{Gauss}schen Tafel die folgenden Werte von $m$:
%CONT>>
